\documentclass[11pt]{article}
\usepackage[paperwidth=7.10in,paperheight=7.75in,left=.43in,right=.43in,top=.28in,bottom=.42in,footskip=.19in]{geometry}
\usepackage{amsmath}
\usepackage[hidelinks]{hyperref}
\pdfinfoomitdate=1
\pdftrailerid{}
\hypersetup{pdfauthor={Kornstein Lab for Antisemitism Research},pdfsubject={Research proposal; illustrative methods, no empirical findings}}
\setlength{\parindent}{0pt}
\setlength{\parskip}{3pt}
\setlength{\abovedisplayskip}{7pt}
\setlength{\belowdisplayskip}{7pt}
\pagestyle{plain}
\makeatletter
\renewcommand\section{\@startsection{section}{1}{0pt}{9pt plus 1pt minus 1pt}{3pt}{\normalfont\large\bfseries}}
\makeatother
\newcommand{\papertitle}[1]{%
  \hypersetup{pdftitle={#1}}%
  \begin{center}
    {\fontsize{18.5}{22}\selectfont #1\par}
    \vspace{6pt}
    {\fontsize{13}{16}\selectfont Kornstein Lab for Antisemitism Research\par}
    {\fontsize{12}{15}\selectfont Research proposal\par}
  \end{center}
  \vspace{2pt}
}
\newenvironment{paperabstract}{\begin{list}{}{\leftmargin=21pt\rightmargin=21pt}\item[]\begin{center}\textbf{Abstract}\end{center}\vspace{-3pt}\small}{\end{list}\vspace{2pt}}
\newenvironment{methodlist}{\begin{enumerate}\setlength{\itemsep}{1pt}\setlength{\parsep}{0pt}\setlength{\topsep}{2pt}}{\end{enumerate}}

\begin{document}
\papertitle{Evaluating algorithms for antisemitism detection}
\begin{paperabstract}
This proposed study evaluates how algorithms distinguish antisemitic expression from quotation, criticism, and other contextual uses. It examines classification errors, consistency, and performance across different kinds of text.
\end{paperabstract}

\section{Introduction}
We propose a contextual evaluation of algorithms for detecting antisemitism. The aim is to compare model outputs with reviewed annotations and examine failure modes across linguistic settings.

\section{Evaluation framework}
An illustrative classification measure is
\[
  F_1=\frac{2PR}{P+R},\qquad
  P=\frac{\mathrm{TP}}{\mathrm{TP}+\mathrm{FP}},\quad
  R=\frac{\mathrm{TP}}{\mathrm{TP}+\mathrm{FN}}.
\]
Precision $P$ and recall $R$ summarize different error costs. We would report both with uncertainty and contextual error analysis, rather than rely on a single score. Undefined ratios require explicit handling.

\section{Proposed methodology}
\begin{methodlist}
  \item Contextual annotation.
  \item Algorithm benchmarking.
  \item Error and robustness analysis.
\end{methodlist}

\section{Limitations}
Annotation judgments and corpus selection may limit transfer across languages and contexts. Model performance remains to be measured; this proposal reports no benchmark results.
\end{document}
